\documentclass[12pt]{article}

% Idioma e codificação
\usepackage[brazil]{babel}
\usepackage[utf8]{inputenc}
\usepackage[T1]{fontenc}

% Matemática, listas e tabelas
\usepackage{amsmath, amssymb}
\usepackage{enumitem}
\usepackage{float} 
\usepackage{tabularx}
\usepackage{booktabs}
\usepackage{colortbl}

% Layout e cores
\usepackage{geometry}
\usepackage{graphicx}
\usepackage{xcolor}

% Cores institucionais
\definecolor{maincolor}{RGB}{0,100,50}
\definecolor{forestgreen}{RGB}{0,100,50}
\definecolor{tableheader}{RGB}{230,240,230}

% Cabeçalho e rodapé
\usepackage{fancyhdr}

% Blocos
\usepackage[most]{tcolorbox}
\newtcolorbox{block}[1]{%
  colback=green!5,
  colframe=maincolor,
  coltitle=white,
  title=#1,
  fonttitle=\bfseries,
  boxrule=1.5pt,
  arc=3mm,
  left=6pt,
  right=6pt,
  top=6pt,
  bottom=6pt
}

% Margens
\geometry{margin=2.5cm}

% Fancyhdr
\pagestyle{fancy}
\fancyhf{}
\renewcommand{\headrulewidth}{2pt}
\renewcommand{\footrulewidth}{2pt}
\renewcommand{\headrule}{\hbox to\headwidth{\color{forestgreen}\leaders\hrule height \headrulewidth\hfill}}
\renewcommand{\footrule}{\hbox to\headwidth{\color{forestgreen}\leaders\hrule height \footrulewidth\hfill}}
\setlength{\headheight}{52pt}
\addtolength{\topmargin}{-20pt}

\fancyhead[L]{%
  \includegraphics[height=3.2cm]{rv.jpeg}
}

\fancyhead[R]{%
  \raggedleft
  \textcolor{forestgreen}{%
    \small
    \textbf{Lista de Exercícios}\\
    Me.~Sergio Souza Novak\\
    \textit{Organização de Pastas Windows}
    }
}
    
\fancyfoot[C]{\textcolor{forestgreen}{\thepage}}
    
\begin{document}
    
\vspace*{0.2cm}

\section*{Exercício 01 – Empresa de Sementes}

Organize a estrutura de pastas do sistema de arquivos da empresa de sementes conforme a hierarquia abaixo. Cada pasta representa um departamento ou categoria de informações.

\begin{table}[H]
\centering
\small
\textbf{Estrutura de Pastas Sugerida}\\
\begin{tabular}{l}
\hline
\rowcolor{tableheader} \textbf{Sementes} \\ \hline
\quad Produtos \\
\quad\quad Milho \\
\quad\quad Soja \\
\quad\quad Trigo \\
\quad\quad Feijão \\
\quad\quad Sementes\_Híbridas \\
\quad Fornecedores \\
\quad Pedidos \\
\quad\quad 2025 \\
\quad\quad\quad T1 \\
\quad\quad\quad T2 \\
\quad\quad 2026 \\
\quad\quad\quad T1 \\
\quad Clientes \\
\quad\quad Norte \\
\quad\quad Sul \\
\quad\quad Leste \\
\quad\quad Oeste \\
\quad Marketing \\
\quad\quad Campanhas \\
\quad\quad RedesSociais \\
\quad\quad Email\_Marketing \\
\hline
\end{tabular}
\end{table}

\begin{block}{Instruções}
Crie todas as pastas acima no diretório \texttt{C:/Documents/Sementes}. Certifique-se de que a hierarquia está correta.
\end{block}
\vspace*{5cm}

\section*{Exercício 02 – Empresa Agropecuária}

Organize os dados da fazenda considerando os diferentes setores de produção, animais e funcionários.

\begin{table}[H]
\centering
\small
\textbf{Estrutura de Pastas Sugerida}\\
\begin{tabular}{l}
\hline
\rowcolor{tableheader} \textbf{Agropecuaria} \\ \hline
\quad Vacas \\
\quad\quad Fazenda\_Sul \\
\quad\quad\quad Leiteiras \\
\quad\quad\quad Bezerros \\
\quad\quad Fazenda\_Norte \\
\quad\quad\quad Leiteiras \\
\quad\quad\quad Bezerros \\
\quad Ovinos \\
\quad Suinos \\
\quad Funcionarios \\
\quad\quad Ordenhadores \\
\quad\quad\quad Manha \\
\quad\quad\quad Tarde \\
\quad\quad Veterinarios \\
\quad\quad\quad Manha \\
\quad\quad Supervisores \\
\quad Relatorios \\
\quad\quad Producao \\
\quad\quad SaudeAnimal \\
\quad\quad ManutencaoEquipamentos \\
\quad Equipamentos \\
\quad Financeiro \\
\hline
\end{tabular}
\end{table}

\begin{block}{Instruções}
Crie cada pasta e subpasta conforme indicado, garantindo consistência na hierarquia.
\end{block}
\vspace*{5cm}

\section*{Exercício 03 – Clínica Veterinária}

Organize os arquivos da clínica considerando pacientes, atendimentos, vacinas, exames e faturamento.

\begin{table}[H]
\centering
\small
\textbf{Estrutura de Pastas Sugerida}\\
\begin{tabular}{l}
\hline
\rowcolor{tableheader} \textbf{ClinicaVeterinaria} \\ \hline
\quad Pacientes \\
\quad\quad Cachorros \\
\quad\quad\quad 2025 \\
\quad\quad\quad 2026 \\
\quad\quad Gatos \\
\quad\quad Aves \\
\quad\quad Outros \\
\quad Atendimentos \\
\quad\quad Janeiro \\
\quad\quad Fevereiro \\
\quad\quad ... \\
\quad Vacinas \\
\quad\quad Antirrabica \\
\quad\quad V10 \\
\quad\quad GripeAnimal \\
\quad Exames \\
\quad\quad Sanguineo \\
\quad\quad Urina \\
\quad\quad Imagem \\
\quad Faturamento \\
\quad Estoque \\
\hline
\end{tabular}
\end{table}

\begin{block}{Instruções}
Implemente a estrutura no Windows, criando pastas e subpastas exatamente como na tabela acima.
\end{block}
\vspace*{5cm}

\section*{Exercício 04 – Petshop}

Organize o sistema de arquivos da loja de animais considerando estoque, vendas, clientes e promoções.

\begin{table}[H]
\centering
\small
\textbf{Estrutura de Pastas Sugerida}\\
\begin{tabular}{l}
\hline
\rowcolor{tableheader} \textbf{Petshop} \\ \hline
\quad Estoque \\
\quad\quad Racoes \\
\quad\quad Brinquedos \\
\quad\quad Higiene \\
\quad\quad Medicamentos \\
\quad Vendas \\
\quad\quad T1\_2026 \\
\quad\quad\quad Dinheiro \\
\quad\quad\quad Cartao \\
\quad Clientes \\
\quad\quad PF \\
\quad\quad\quad Norte \\
\quad\quad\quad Sul \\
\quad\quad PJ \\
\quad\quad\quad Norte \\
\quad Promocoes \\
\quad Funcionarios \\
\hline
\end{tabular}
\end{table}

\begin{block}{Instruções}
Crie as pastas no diretório \texttt{C:/Documents/Petshop} e verifique se a hierarquia está correta.
\end{block}

\end{document}